
%-------------------------------------------------%
% 2025-04-01 - Emerson Ribeiro de Mello - mello@ifsc.edu.br
%-------------------------------------------------%
\documentclass[aspectratio=169]{beamer}

\usepackage[english,brazil]{babel}
\usepackage[style=abnt,noslsn,justify]{biblatex}
\renewcommand*{\bibfont}{\small}
\usepackage{csquotes}
\addbibresource{referencias_trabalho.bib}

\definecolor[named]{novaCor}{HTML}{303A52}
\definecolor{verdeescuro}{RGB}{0,112,19}

\usetheme{ifsclean}

% -------------------------------------------------%
%              Título 
% -------------------------------------------------%
\title{OWASP WebGoat e WebWolf}
\subtitle{Laboratório de Segurança em Aplicações Web}
\author{
  Giuliano Marcus Bianco \\[3pt]
  Luis Renato Freitas de Almeida \\[3pt]
  Nicolas Arthur Raulino Oliveira
}
\institute{Instituto Federal de Santa Catarina}
\date{Novembro de 2025}

% -------------------------------------------------%
\begin{document}

\begin{frame}
    \titlepage
\end{frame}

% -------------------------------------------------%
\section{Introdução}

\begin{frame}{Contextualização}
    \begin{itemize}
        \item O crescimento das aplicações web trouxe novos desafios de segurança.
        \item A OWASP promove boas práticas e ferramentas para mitigar vulnerabilidades.
        \item WebGoat e WebWolf são plataformas didáticas para aprendizado prático de segurança.
    \end{itemize}
\end{frame}

% -------------------------------------------------%
\section{A OWASP}

\begin{frame}{O que é a OWASP}
    \begin{itemize}
        \item Comunidade global dedicada à segurança de software.
        \item Produz guias, ferramentas e projetos abertos, como o \textbf{OWASP Top 10}.
        \item O objetivo é ajudar desenvolvedores a criar aplicações mais seguras.
    \end{itemize}
\end{frame}

% -------------------------------------------------%
\section{OWASP WebGoat}

\begin{frame}{O que é o WebGoat}
    \begin{itemize}
        \item Aplicação web vulnerável criada intencionalmente.
        \item Desenvolvida em Java e executada no servidor Apache Tomcat.
        \item Usada para ensinar como explorar e corrigir falhas de segurança.
    \end{itemize}
\end{frame}

\begin{frame}{Exemplos de vulnerabilidades abordadas}
    \begin{itemize}
        \item Injeção de SQL (\textit{SQL Injection});
        \item Quebra de autenticação e controle de sessão;
        \item Cross-Site Scripting (XSS);
        \item Exposição de dados sensíveis;
        \item Controle de acesso incorreto.
    \end{itemize}
\end{frame}

% -------------------------------------------------%
\section{OWASP WebWolf}

\begin{frame}{O que é o WebWolf}
    \begin{itemize}
        \item Ferramenta complementar ao WebGoat.
        \item Simula o ambiente de um atacante externo.
        \item Permite receber e-mails falsos, capturar dados e testar exploits.
    \end{itemize}
\end{frame}

\begin{frame}{Integração com o WebGoat}
    \begin{itemize}
        \item WebGoat envia dados ao WebWolf durante atividades práticas.
        \item Juntos, simulam cenários reais de ataque e defesa.
        \item Ideal para treinamento ético em segurança web.
    \end{itemize}
\end{frame}

% -------------------------------------------------%
\section{Demonstração}

\begin{frame}{Ambiente de Teste}
    \begin{itemize}
        \item Instalação via Docker Compose com dois containers: \texttt{webgoat} e \texttt{webwolf}.
        \item Endereços padrão:
              \begin{itemize}
                  \item WebGoat: \texttt{http://localhost:8080/WebGoat}
                  \item WebWolf: \texttt{http://localhost:9090/WebWolf}
              \end{itemize}
        \item Aulas interativas e relatórios automáticos.
    \end{itemize}
\end{frame}

\begin{frame}{Exemplo de Exercício}
    \begin{itemize}
        \item Exploração de um formulário vulnerável (ex: login).
        \item Observação dos logs e e-mails capturados no WebWolf.
        \item Discussão sobre a mitigação da falha.
    \end{itemize}
\end{frame}

% -------------------------------------------------%
\section{Conclusão}

\begin{frame}{Considerações Finais}
    \begin{itemize}
        \item WebGoat e WebWolf permitem aprendizado prático e seguro de vulnerabilidades.
        \item Ajudam a compreender como proteger aplicações contra ataques reais.
        \item São ferramentas essenciais para o ensino de segurança da informação.
    \end{itemize}
\end{frame}

% -------------------------------------------------%
\begin{frame}[allowframebreaks]{Referências}
    \nocite{*}
    \printbibliography
\end{frame}

\end{document}
